\documentclass[11pt,a4paper]{article}

\usepackage[utf8]{inputenc}
\usepackage[T1]{fontenc}
\usepackage[brazil]{babel}
\usepackage{amsmath,amssymb,amsfonts,amsthm}
\usepackage{geometry}
\usepackage{booktabs}
\usepackage{xcolor}
\usepackage{tcolorbox}
\usepackage{enumitem}
\usepackage{tabularx}
\usepackage{tikz}
\usetikzlibrary{arrows.meta,trees,positioning}
\usepackage{hyperref}

\geometry{
  top=2cm,
  bottom=2cm,
  left=2cm,
  right=2cm
}

\setlist{itemsep=2pt, topsep=2pt}

\hypersetup{
  colorlinks=true,
  linkcolor=blue!70!black,
  citecolor=blue!70!black,
  urlcolor=blue!70!black
}

\title{\vspace{-1cm}\textbf{Exercício 4 --- Refutação por Resolução}\\
\large Inteligência Artificial --- Lógica de Predicados de Primeira Ordem}
\author{\textbf{André Moreno Goveia} \\ NUSP: 13682785}
\date{\today}

\begin{document}

\maketitle

\vspace{-0.4cm}
\hrule
\vspace{0.4cm}

\section*{Contextualização do Problema}

Considere o seguinte conteúdo de uma base de conhecimento ($KB$):
\begin{enumerate}[label=\textbf{\arabic*.}]
    \item \textit{Todo engenheiro que trabalha na Petrobrás terá estabilidade no emprego nos próximos 5 anos e será reconhecido nacionalmente por seu trabalho nos próximos 10 anos.}
    \item \textit{Existe um engenheiro que trabalha na Petrobrás e será presidente do STF.}
\end{enumerate}

% ==============================================================================
% ITEM A
% ==============================================================================
\section*{Item (a) --- Representação em Lógica de Predicados}

\subsection*{1. Vocabulário, Predicados, Funções e Variáveis}

\begin{itemize}
    \item \textbf{Domínio do Discurso ($\mathcal{D}$):} Conjunto de pessoas (indivíduos).
    \item \textbf{Predicados:}
    \begin{itemize}
        \item $E(x)$: ``$x$ é engenheiro''
        \item $P(x)$: ``$x$ trabalha na Petrobrás''
        \item $S(x)$: ``$x$ terá estabilidade no emprego nos próximos 5 anos''
        \item $R(x)$: ``$x$ será reconhecido nacionalmente por seu trabalho nos próximos 10 anos''
        \item $J(x)$: ``$x$ será presidente do STF''
    \end{itemize}
    \item \textbf{Variáveis:} $x, y \in \mathcal{D}$ (representam indivíduos arbitrários do domínio).
    \item \textbf{Constantes e Funções:} Não há constantes nem funções primitivas na especificação inicial. A constante de Skolem $a$ surge da eliminação de quantificadores existenciais na etapa de conversão clausal.
\end{itemize}

\subsection*{2. Formalização das Sentenças da Base}

\begin{itemize}
    \item \textbf{Sentença 1 ($S_1$):}
    \[
    \forall x \, \Big( (E(x) \land P(x)) \to (S(x) \land R(x)) \Big)
    \]
    \item \textbf{Sentença 2 ($S_2$):}
    \[
    \exists x \, \Big( E(x) \land P(x) \land J(x) \Big)
    \]
\end{itemize}

\subsection*{3. Conversão da Base ($KB$) para a Forma Normal Clausal (CNF)}

\paragraph{Transformação de $S_1$:}
\begin{enumerate}[label=\alph*)]
    \item Eliminação da implicação ($A \to B \equiv \neg A \lor B$):
    $\forall x \, \big( \neg(E(x) \land P(x)) \lor (S(x) \land R(x)) \big)$
    \item Aplicação de De Morgan:
    $\forall x \, \big( \neg E(x) \lor \neg P(x) \lor (S(x) \land R(x)) \big)$
    \item Distributividade de $\lor$ sobre $\land$:
    $\forall x \, \big( (\neg E(x) \lor \neg P(x) \lor S(x)) \land (\neg E(x) \lor \neg P(x) \lor R(x)) \big)$
    \item Eliminação do quantificador universal $\forall$, obtendo duas cláusulas disjuntivas:
    \begin{align*}
        C_1: & \quad \neg E(x) \lor \neg P(x) \lor S(x) \\
        C_2: & \quad \neg E(x) \lor \neg P(x) \lor R(x)
    \end{align*}
\end{enumerate}

\paragraph{Transformação de $S_2$:}
\begin{enumerate}[label=\alph*)]
    \item Skolemização de $\exists x$ (introdução da constante de Skolem $a$):
    $E(a) \land P(a) \land J(a)$
    \item Separação em cláusulas unitárias:
    \begin{align*}
        C_3: & \quad E(a) \\
        C_4: & \quad P(a) \\
        C_5: & \quad J(a)
    \end{align*}
\end{enumerate}

Portanto, o conjunto de cláusulas da base é:
\[
KB = \{ C_1, C_2, C_3, C_4, C_5 \}
\]

% ==============================================================================
% ITEM B
% ==============================================================================
\section*{Item (b) --- Análise da Sentença 3 por Refutação}

Deseja-se verificar se a seguinte sentença decorre logicamente da base:
\begin{enumerate}
    \item[\textbf{3.}] \textit{Todo engenheiro que trabalha na Petrobrás será reconhecido nacionalmente por seu trabalho nos próximos 10 anos e será presidente do STF.}
\end{enumerate}

\subsection*{1. Formalização da Meta ($\alpha$)}
\[
\alpha: \forall x \, \Big( (E(x) \land P(x)) \to (R(x) \land J(x)) \Big)
\]

\subsection*{2. Negação da Meta ($\neg \alpha$) e Clausalização}
Pelo método de refutação por resolução, adiciona-se a negação da meta ($\neg \alpha$) a $KB$:
\begin{align*}
\neg \alpha & \equiv \neg \Big( \forall x \, \big( (E(x) \land P(x)) \to (R(x) \land J(x)) \big) \Big) \\
& \equiv \exists x \, \neg \Big( \neg(E(x) \land P(x)) \lor (R(x) \land J(x)) \Big) \\
& \equiv \exists x \, \Big( E(x) \land P(x) \land \neg(R(x) \land J(x)) \Big) \\
& \equiv \exists x \, \Big( E(x) \land P(x) \land (\neg R(x) \lor \neg J(x)) \Big)
\end{align*}

Ao eliminar o quantificador existencial $\exists x$, devemos introduzir uma \textbf{nova constante de Skolem}, $b$ (com $b \neq a$), pois a Skolemização não permite reutilizar uma constante pré-existente sem perda de consistência semântica:
\begin{align*}
C_6: & \quad E(b) \\
C_7: & \quad P(b) \\
C_8: & \quad \neg R(b) \lor \neg J(b)
\end{align*}

\subsection*{3. Aplicação das Regras de Resolução}
Buscamos derivar a cláusula vazia ($\square$) a partir de $KB \cup \{C_6, C_7, C_8\}$:
\begin{enumerate}[label=\arabic*.]
    \item Resolução de $C_2$ com $C_6$ (MGU $\theta_1 = \{x / b\}$):
    \[
    \mathrm{Res}(C_2, C_6) = \neg P(b) \lor R(b) \quad (C_9)
    \]
    \item Resolução de $C_9$ com $C_7$:
    \[
    \mathrm{Res}(C_9, C_7) = R(b) \quad (C_{10})
    \]
    \item Resolução de $C_8$ com $C_{10}$:
    \[
    \mathrm{Res}(C_8, C_{10}) = \neg J(b) \quad (C_{11})
    \]
\end{enumerate}

Para obter a cláusula vazia ($\square$), o literal negativo $\neg J(b)$ precisaria resolver com o literal positivo $J(a)$ de $C_5$. Todavia:
\[
\mathrm{MGU}(J(b), J(a)) = \text{Falha (não unificam, pois $a$ e $b$ são constantes distintas)}.
\]
Nenhuma outra inferência gera literais complementares adicionais. O conjunto de cláusulas é estritamente consistente.

\begin{tcolorbox}[colback=red!5!white,colframe=red!70!black,title=\textbf{Conclusão do Item (b)}]
\textbf{A sentença 3 NÃO é consequência lógica da base de conhecimento ($KB \not\models \alpha$).} \\
Não é possível derivar a contradição ($\square$). De fato, a base afirma apenas que \emph{existe algum} engenheiro que presidirá o STF, o que não autoriza concluir que \emph{todos} os engenheiros da Petrobrás alcançarão essa posição.
\end{tcolorbox}

% ==============================================================================
% ITEM C
% ==============================================================================
\section*{Item (c) --- Demonstração da Sentença 4 por Refutação}

Deseja-se verificar se a seguinte sentença decorre logicamente da base:
\begin{enumerate}
    \item[\textbf{4.}] \textit{Algum engenheiro que trabalha na Petrobrás será reconhecido nacionalmente por seu trabalho nos próximos 10 anos e será presidente do STF.}
\end{enumerate}

\subsection*{1. Formalização da Meta ($\beta$)}
\[
\beta: \exists x \, \Big( E(x) \land P(x) \land R(x) \land J(x) \Big)
\]

\subsection*{2. Negação da Meta ($\neg \beta$) e Clausalização}
\begin{align*}
\neg \beta & \equiv \neg \Big( \exists x \, (E(x) \land P(x) \land R(x) \land J(x)) \Big) \\
& \equiv \forall x \, \Big( \neg E(x) \lor \neg P(x) \lor \neg R(x) \lor \neg J(x) \Big)
\end{align*}

Como a negação contém apenas um quantificador universal, na forma clausal o símbolo quantificado torna-se uma variável universalmente quantificada (variável livre). Padronizando com a variável $y$:
\[
C_6': \quad \neg E(y) \lor \neg P(y) \lor \neg R(y) \lor \neg J(y)
\]

\subsection*{3. Tabela de Resolução Passo a Passo}

Aplicando o princípio da resolução ao conjunto $KB \cup \{C_6'\}$:

\begin{center}
\small
\resizebox{\linewidth}{!}{%
\begin{tabular}{c l l c l}
\toprule
\textbf{Passo} & \textbf{Cláusula Pai 1} & \textbf{Cláusula Pai 2} & \textbf{MGU ($\theta$)} & \textbf{Resolvente Derivado} \\
\midrule
1 & $C_6'$ & $C_5: J(a)$ & $\{y / a\}$ & $C_7': \neg E(a) \lor \neg P(a) \lor \neg R(a)$ \\
2 & $C_7'$ & $C_4: P(a)$ & $\emptyset$ & $C_8': \neg E(a) \lor \neg R(a)$ \\
3 & $C_8'$ & $C_3: E(a)$ & $\emptyset$ & $C_9': \neg R(a)$ \\
4 & $C_2: \neg E(x) \lor \neg P(x) \lor R(x)$ & $C_9'$ & $\{x / a\}$ & $C_{10}': \neg E(a) \lor \neg P(a)$ \\
5 & $C_{10}'$ & $C_4: P(a)$ & $\emptyset$ & $C_{11}': \neg E(a)$ \\
6 & $C_{11}'$ & $C_3: E(a)$ & $\emptyset$ & $C_{12}': \mathbf{\square} \quad \text{\textbf{(Cláusula Vazia)}}$ \\
\bottomrule
\end{tabular}%
}
\end{center}

\subsection*{4. Árvore de Resolução por Refutação}

\begin{center}
\resizebox{0.88\linewidth}{!}{%
\begin{tikzpicture}[
  level distance=1.25cm,
  sibling distance=4.2cm,
  every node/.style={draw, rounded corners, fill=blue!5, font=\small, align=center, inner sep=4pt},
  edge from parent/.style={draw, -{Latex[length=2mm]}}
]
\node[fill=red!15, very thick] {$\mathbf{\square}$ \textbf{(Contradição)}}
  child {
    node {$C_3: E(a)$}
  }
  child {
    node {$C_{11}': \neg E(a)$}
    child {
      node {$C_4: P(a)$}
    }
    child {
      node {$C_{10}': \neg E(a) \lor \neg P(a)$}
      child {
        node {$C_2: \neg E(x) \lor \neg P(x) \lor R(x)$}
      }
      child {
        node {$C_9': \neg R(a)$}
        child {
          node {$C_3: E(a)$}
        }
        child {
          node {$C_8': \neg E(a) \lor \neg R(a)$}
          child {
            node {$C_4: P(a)$}
          }
          child {
            node {$C_7': \neg E(a) \lor \neg P(a) \lor \neg R(a)$}
            child {
              node {$C_5: J(a)$}
            }
            child {
              node {$C_6': \neg E(y) \lor \neg P(y) \lor \neg R(y) \lor \neg J(y)$}
            }
          }
        }
      }
    }
  };
\end{tikzpicture}%
}
\end{center}

\begin{tcolorbox}[colback=green!5!white,colframe=green!60!black,title=\textbf{Conclusão do Item (c)}]
\textbf{A sentença 4 É consequência lógica da base de conhecimento ($KB \models \beta$).} \\
A cláusula vazia ($\square$) foi formalmente deduzida a partir de $KB \cup \{\neg \beta\}$. Como a refutação gerou contradição, a hipótese de falsidade de $\beta$ é insustentável. O engenheiro da Petrobrás cuja existência no STF é garantida pela Sentença 2 necessariamente usufrui de reconhecimento nacional em virtude da Sentença 1.
\end{tcolorbox}

\end{document}
